% SPDX-License-Identifier: LPPL-1.3c
% Copyright (C) 2026 Christophe Jorssen
% Author and Current Maintainer: Christophe Jorssen <christophe.jorssen@gmail.com>
% LPPL maintenance status: maintained
% This file is part of luafemm. See LICENSE and LICENSES.md for its terms.
\section{Static models and tagged frontiers}
\label{sec:static}
This chapter extends the planar, direct-current model: the computational domain
can follow a TikZ path; named boundaries can be paired; materials can have a
linear permeability tensor; permanent magnets can have a spatially varying
direction; and named windings can supply their current and report flux linkage.
All features use the same pure-Lua solver and generic TeX interface.
They do not introduce an axisymmetric or alternating-current formulation.

\subsection{Decide what each path means}
A magnetic drawing can contain four different kinds of path. A |femm/region|
path assigns material. A |femm/domain| path determines where the equation is
solved. A |femm/frontier| path assigns a physical boundary condition. A
|femm/profile| or |femm/selection| path measures an existing solution. Ordinary
TikZ paths only draw. Keeping these roles explicit makes the diagram easier to
edit: moving a profile cannot accidentally move a physical boundary.

Start with a coarse |mesh size| and modest |curve tolerance|. Finish the geometry,
sources and styles, then reduce both tolerances where the final result requires
it. Decreasing the curve tolerance improves the polygonal representation;
decreasing the mesh size improves the finite-element solution on that
representation. They are separate convergence checks.

The order of operations is:
\begin{enumerate}
\item Define material styles and reusable boundary/circuit styles.
\item Start |tikzpicture| with |femm/problem|. Circuit and pairing declarations in
  picture options are queued until the model exists, just like |femm/boundary|.
\item Capture the computational domain, material regions and tagged paths.
  Their declaration order is free, except that material paint order still matters.
  Declare each circuit before the mesh is generated; its paths may come first.
\item Mesh or solve. A field pic, profile sample or measurement may trigger this
  step automatically. Finish \emph{all} physical declarations first.
\item Draw the field, take measurements and plot their named snapshots.
\end{enumerate}
Physical declarations persist beyond the graphical scope that captured them.
They cannot be added after meshing. Reusable |.style|, |.append style| and
parameterised PGF styles work throughout this chapter. Numerical keys accept
PGF expressions; values passed to Lua are validated finite numbers.

\subsection{A computational domain drawn with TikZ}
\begin{codeexample}[width=6cm]
\begin{tikzpicture}[femm/problem={
  xmin=-22,xmax=22,ymin=-17,ymax=17,
  mesh size=4}]
  \path[femm/domain]
    (-18,-12)--(18,-8)--(14,13)--(-14,10)--cycle;
  \draw[thick,femm/frontier={name=outer,
    y coefficient=.2}]
    (-18,-12)--(18,-8)--(14,13)--(-14,10)--cycle;
  \pic[femm/lines=7] {femm field};
\end{tikzpicture}
\end{codeexample}

\begin{key}{/tikz/femm/domain=\marg{options} (default \normalfont empty)}
Capture one closed, possibly compound path as the computational domain. At most
one domain declaration is allowed per problem. The TikZ fill rule applies:
|even odd rule| conveniently describes holes. Holes in this domain exclude
triangles; holes in a material region merely expose the previous material.
The problem rectangle is still required as a bounding/staging box, and must
contain every physical path. Its sides cease to be physical boundaries when a
custom domain is present. Unlabelled custom-domain boundaries default to
$A_z=0$, independently of rectangle-side declarations.

With no |femm/domain|, the original rectangular domain and its side conditions
remain in effect. Path frontiers, custom domains and circuits require
|mesher=delaunay| and |field model=open|. They do not modify an imported FEM
model or an ideal-circuit model.
\end{key}
\begin{key}{/femm/domain/curve tolerance=\meta{distance} (initially .03)}
Positive flattening tolerance in model units. One closed domain can contain
several disjoint components. Ordinary TikZ curves, arcs and rounded corners are
flattened after their transformations have been evaluated.
\end{key}

\subsection{Tagging a physical boundary}
\begin{key}{/tikz/femm/frontier=\marg{options} (default \normalfont empty)}
Attach a named condition to the captured path. The path may be open or closed,
but must belong to the graph of physical boundaries/interfaces. An internal
Dirichlet contour or an internal path joining existing boundaries is supported
when it bounds valid faces. An isolated open internal constraint is rejected.
Neumann, Robin and periodic conditions require exterior edges.

The condition follows every split segment during arrangement construction and
meshing. It overrides the default condition on the same exterior edge.
Conflicting explicit tags on coincident edges are rejected. The existing key
|femm/boundary| retains its original meaning: it selects rectangle sides and
\emph{does not} tag the path on which it occurs.
\end{key}
\begin{key}{/femm/frontier/name=\meta{name} (initially empty)}
Required problem-local name. Names are unique within a problem; repeated names
are errors. Names need not be TeX control sequences.
\end{key}
\begin{key}{/femm/frontier/type=\meta{choice} (initially dirichlet)}
An |.is choice| key accepting |dirichlet|, |neumann|, |robin|, |periodic| and
|open|. With outward normal $\mathbf n$, define the positive tangent
$\mathbf t=(-n_y,n_x)$. Dirichlet prescribes $A_z=g$; Neumann prescribes
$\mathbf H\cdot\mathbf t=g$; Robin prescribes
$\mathbf H\cdot\mathbf t=cA_z+g$, with $c\geq0$.
|periodic| marks a path for a separate pairing declaration, whose sign can be
periodic or antiperiodic. |open| is the circular multipole approximation below.
\end{key}
\begin{key}{/femm/frontier/value=\meta{number} (initially 0)}
Constant term of $g=\texttt{value}+\texttt{x coefficient}\,x+
\texttt{y coefficient}\,y$, where $x,y$ are in metres. Potential data use T\,m;
Neumann/Robin data use A/m. Geometry remains in model units.
\end{key}
\begin{key}{/femm/frontier/x coefficient=\meta{number} (initially 0)}
Coefficient of the SI abscissa in $g$.
\end{key}
\begin{key}{/femm/frontier/y coefficient=\meta{number} (initially 0)}
Coefficient of the SI ordinate in $g$.
\end{key}
\begin{key}{/femm/frontier/coefficient=\meta{number} (initially 0)}
Nonnegative Robin coefficient $c$, in A/(T\,m$^2$); other types require zero.
\end{key}
\begin{key}{/femm/frontier/curve tolerance=\meta{distance} (initially .03)}
Positive path-flattening tolerance in model units. Use the same tolerance when
separate paths describe the same curved interface.
\end{key}
\begin{stylekey}{/femm/frontier/field tangent}
Homogeneous Dirichlet symmetry: $A_z=0$ and $B_n=0$.
\end{stylekey}
\begin{stylekey}{/femm/frontier/field normal}
Homogeneous natural symmetry: $H_t=0$. In isotropic, source-free air this means
$B_t=0$. In an anisotropic medium, $\mathbf H$ and $\mathbf B$ need not be parallel;
use the defining condition rather than inferring a direction for $\mathbf B$.
\end{stylekey}

\subsection{Working on a periodic sector}
An associated pair identifies the nodal potentials on two exterior traces.
Both traces are first synchronised by inserting the union of their transformed
vertices into their boundary triangles. The solver eliminates the constrained
unknowns exactly; it does not add a penalty parameter. A reversed sign is carried
through the residual, tangent, nodal updates and cache restoration.

The following sector contains an angularly magnetised permanent magnet. Its
$60^\circ$ pairing is antiperiodic. The angle law is chosen consistently with
that symmetry: the magnetisation rotates by $240^\circ$, namely the spatial
rotation plus a $180^\circ$ reversal.
\luafemmexample{static-frontiers}

\begin{key}{/tikz/femm/pair=\marg{options}}
Declare a pair through ordinary picture, scope or body keys. Picture declarations
are queued; body declarations apply immediately. The referenced paths must exist
by the time meshing starts, and each path belongs to exactly one pair.
\end{key}
\begin{command}{\femmpair\oarg{options}}
Equivalent immediate declaration in a started magnetic picture. Subkeys belong
to |/femm/pair|. Each declaration resets its options before applying the supplied
keys/styles. Multiple pairs may share corner vertices; contradictory nonzero
Dirichlet values at an antiperiodic corner are rejected.
\end{command}
\begin{key}{/femm/pair/first=\meta{name} (initially empty)}
Name of the source trace, declared with |type=periodic|.
\end{key}
\begin{key}{/femm/pair/second=\meta{name} (initially empty)}
Name of the target trace. It must be a rigid image of the first trace.
\end{key}
\begin{key}{/femm/pair/type=\meta{choice} (initially periodic)}
|periodic| imposes $A_2=A_1$; |antiperiodic| imposes $A_2=-A_1$.
The transform maps the first trace onto the second.
\end{key}
\begin{key}{/femm/pair/rotation=\meta{angle} (initially 0)}
Counterclockwise rotation in degrees about the specified centre, followed by
translation. Reflections, scaling and sliding interfaces are not included.
\end{key}
\begin{key}{/femm/pair/center x=\meta{coordinate} (initially 0)}
Abscissa of the rotation centre, in model units.
\end{key}
\begin{key}{/femm/pair/center y=\meta{coordinate} (initially 0)}
Ordinate of the rotation centre, in model units.
\end{key}
\begin{key}{/femm/pair/shift x=\meta{distance} (initially 0)}
Translation along the model $x$ axis after rotation, in model units.
\end{key}
\begin{key}{/femm/pair/shift y=\meta{distance} (initially 0)}
Translation along the model $y$ axis after rotation, in model units.
\end{key}
\begin{key}{/femm/pair/tolerance=\meta{distance} (initially .001)}
Positive matching tolerance in model units. It accommodates PGF coordinate
rounding. Transformed points within this tolerance are projected onto the actual
target trace. Choose a tolerance much smaller than the smallest relevant feature;
it is not a licence to pair physically different boundaries.
\end{key}

Symmetry of the material distribution and sources remains a modelling
responsibility. Synchronisation may reduce the minimum angle near a paired
boundary; inspect the final mesh statistics. No automatic multiplication of
energy, linkage or force measurements is performed. A sector energy can be
summed over identical sectors; force vectors must also be rotated before summing.

\begin{key}{/tikz/femm/repeat=\marg{options}}
Drawing options for the pic |femm repeated field|. The pic reuses solved contours
and the transform of a declared pair. Antiperiodic copies reverse their arrows.
It neither duplicates material regions nor creates new profiles or measurements.
\end{key}
\begin{key}{/femm/repeat/copies=\meta{integer} (initially 2)}
Number of copies including the solved sector, between 1 and 1000.
\end{key}
\begin{key}{/femm/repeat/pair=\meta{integer} (initially 1)}
One-based index of the pairing declaration whose transform is repeated.
\end{key}
For a $60^\circ$ sector, set |femm/repeat={copies=6}| on a
|femm repeated field| pic, as in the complete example above. Draw repeated
outlines with ordinary TikZ paths, without registering additional physical regions.

\subsection{Approximating an open exterior}
On a circle of radius $R$, the decaying exterior angular mode of order $n$ obeys
\[
 A_z\propto r^{-n},\qquad
 \mathbf H\cdot\mathbf t=\frac{n}{\mu_0 R}A_z.
\]
The |open| type installs this nonnegative Robin coefficient. It is exact for
that mode on a true circle in source-free air; on a flattened circle with
several angular modes it is an approximation. Enlarge the circle, refine its
polygonalisation and compare the observables of interest. The native model
checks that the adjacent elements are source-free air and that net longitudinal
current vanishes. Include the return conductor of a winding.

\begin{key}{/femm/frontier/radius=\meta{distance} (initially 1)}
Positive circular radius in model units for |type=open|. The path must be a
complete closed circle about the specified centre; its flattened vertices are
checked allowing curve/coordinate approximation. This key does not draw or
create the domain: declare the circular |femm/domain| separately.
\end{key}
\begin{key}{/femm/frontier/order=\meta{number} (initially 1)}
Positive exterior decay exponent $n$. The default is the first multipole.
Use another value only when the corresponding exterior angular symmetry is
known; this local boundary does not implement a full spectral exterior operator.
\end{key}
\begin{key}{/femm/frontier/center x=\meta{coordinate} (initially 0)}
Circular centre abscissa for |type=open|, in model units.
\end{key}
\begin{key}{/femm/frontier/center y=\meta{coordinate} (initially 0)}
Circular centre ordinate for |type=open|, in model units.
\end{key}

\subsection{Material directions and sheet mixtures}
For a linear material, the solver uses
\[
 \mathbf H=\frac{1}{\mu_0}\mathbf R(\alpha)
 \begin{pmatrix}1/\mu_x&0\\0&1/\mu_y\end{pmatrix}
 \mathbf R(\alpha)^T\mathbf B-H_c\mathbf m.
\]
The full symmetric tensor is used in assembly, field sampling and energy.
Its principal axes and the coercive direction are independent.

\begin{key}{/tikz/femm/materials/mu x=\meta{number} (initially empty)}
Positive relative permeability along the first material principal axis.
An omitted value inherits |mu r| or the selected library material.
\end{key}
\begin{key}{/tikz/femm/materials/mu y=\meta{number} (initially empty)}
Positive relative permeability along the second principal axis. These keys are
for linear media. Unequal principal permeabilities combined with a B--H curve
are rejected; nonlinear anisotropy needs a separate constitutive model.
\end{key}
\begin{key}{/tikz/femm/materials/material angle=\meta{angle} (initially 0)}
Counterclockwise orientation of the first principal axis in the physical model
frame, in degrees. A TikZ transformation of a region changes its geometry; it
does not silently rotate this numerical material direction. Set the angle
explicitly when a component's physical orientation changes. Rotating the entire
picture for presentation preserves the model-relative direction.
\end{key}
\begin{key}{/tikz/femm/materials/lamination=\meta{choice} (initially inherit)}
|inherit| uses a supported catalogue choice or defaults to |none|. The
|parallel x| and |parallel y| choices homogenise linear iron/air
sheets parallel to the selected principal axis, before applying |material angle|.
For fill $f$, permeability parallel to the sheets is $f\mu_r+1-f$; across them
it is $(f/\mu_r+1-f)^{-1}$. Nonlinear sheet mixtures and laminated permanent
magnets are rejected. The |stacked| choice mixes along both field-plane directions; its complex
harmonic extension and thickness parameter are detailed in Section~\ref{sec:losses}.
Catalogue nonlinear curves retain their existing static mixture interpretation.
\end{key}
\begin{key}{/tikz/femm/materials/lamination fill=\meta{fraction} (initially empty)}
Iron fraction $0<f\leq1$. Empty inherits the catalogue value, otherwise defaults
to one. A fraction below one requires a lamination choice.
\end{key}
\begin{key}{/tikz/femm/materials/conductivity=\meta{number} (initially empty)}
Nonnegative conductivity in S/m. An omitted value inherits the catalogue value,
converted from MS/m, or defaults to zero for a custom material. Conductivity
enters the DC winding resistance. It does not introduce eddy currents.
\end{key}

\subsection{Spatial magnetisation without executable expressions}
The region keys below describe the direction of $H_c\mathbf m$. Magnitude and
recoil/B--H law still come from the material. The existing
|magnetization angle| supplies the offset $\alpha_0$ in degrees. Centres and
angles are defined in the physical model frame, independently of drawing styles.

\begin{key}{/femm/region/magnetization=\meta{choice} (initially uniform)}
|uniform| uses the constant offset. |radial| uses
$\alpha=\alpha_0+\theta$; |tangential| uses
$\alpha=\alpha_0+90^\circ+\theta$; |angular| uses
$\alpha=\alpha_0+k\theta$, where
$\theta=\operatorname{atan2}(y-y_c,x-x_c)$ expressed in degrees.
For example |angular,magnetization order=2| describes a cylindrical Halbach
angle law. The family is passive data: arbitrary Lua expressions are not run.
\end{key}
\begin{key}{/femm/region/magnetization center x=\meta{coordinate} (initially 0)}
Abscissa $x_c$, in model units.
\end{key}
\begin{key}{/femm/region/magnetization center y=\meta{coordinate} (initially 0)}
Ordinate $y_c$, in model units.
\end{key}
\begin{key}{/femm/region/magnetization order=\meta{number} (initially 1)}
Multiplier $k$ for |angular|; ignored by the other presets. Positive, zero and
negative finite values are accepted. Noninteger orders introduce an angular
branch discontinuity; use them only on a domain where that is intentional.
\end{key}

Three symmetric quadrature points average the coercive direction in each
triangle. The weak source, reported element $\mathbf H$ and energy use the same
average. Its magnitude can be below one on coarse elements. Refine the magnet
until the angular variation is resolved. A polar law is undefined at its centre;
annular or sector-shaped magnets naturally avoid that point.

The following complete example combines an angularly magnetised annulus, a
circular open boundary and an ordinary field profile through its bore.
\luafemmexample{static-open}

\subsection{Named circuits and signed-turn windings}
A circuit supplies a terminal current or voltage. It contains one or more named
windings, each made from region legs. A leg's signed turn count determines the
sign of its longitudinal current density. Group outgoing and returning legs
under the same winding name. The default winding name is |main|.

\begin{key}{/tikz/femm/circuit=\marg{options}}
Declare a named circuit in picture options, scope options or the picture body.
Picture options are queued until the problem is created. Duplicate circuit
names are errors; every declared circuit must have at least one active leg.
\end{key}
\begin{command}{\femmcircuit\oarg{options}}
Equivalent immediate command in a magnetic picture. Subkeys are in
|/femm/circuit|, and defaults reset for each declaration.
\end{command}
\begin{key}{/femm/circuit/name=\meta{name} (initially empty)}
Required problem-local circuit identifier.
\end{key}
\begin{key}{/femm/circuit/type=\meta{choice} (initially series)}
|series| passes the terminal current through every winding. |parallel| splits
it among named windings according to their DC conductances. Legs of one winding
are always in series. Distinct parallel branches therefore need distinct
|winding| names, even if each branch contains two return-path legs.
\end{key}
\begin{key}{/femm/circuit/drive=\meta{choice} (initially current)}
|current| or |voltage|. This is a static resistive circuit model. There is no
induced voltage, skin effect or time-dependent circuit equation.
\end{key}
\begin{key}{/femm/circuit/current=\meta{number} (initially 0)}
Signed terminal current in amperes for current drive. Leave voltage at zero.
\end{key}
\begin{key}{/femm/circuit/voltage=\meta{number} (initially 0)}
Signed terminal voltage in volts for voltage drive. Leave current at zero.
\end{key}
\begin{key}{/femm/region/circuit=\meta{name} (initially empty)}
Associate a material region with a circuit. Its prescribed |current density|
must be zero; the circuit supplies its excitation. This also applies to
catalogue materials that carry a nonzero default source: explicitly set
|current density=0| for those materials.
\end{key}
\begin{key}{/femm/region/winding=\meta{name} (initially main)}
Group related outgoing/returning legs within the circuit.
\end{key}
\begin{key}{/femm/region/turns=\meta{number} (initially 1)}
Finite, nonzero signed turn count $N$. Negative turns reverse $J_z$.
A fractional effective count is allowed for homogenised winding models.
\end{key}
\begin{key}{/femm/region/winding fill=\meta{fraction} (initially 1)}
Conductor area fraction $0<f\leq1$. It changes resistance, not the total
ampere-turns or the material permeability.
\end{key}
\begin{key}{/femm/region/conductor length=\meta{distance} (initially 0)}
Electrical length per turn of this leg, in model units. Zero inherits problem
|depth|. An explicit larger value can account for allocated end-connection
resistance; those extra lengths do not create an out-of-plane magnetic model.
\end{key}

Let $S$ be a leg's \emph{active} area after all material paint operations, $\ell$
its electrical length, $\sigma$ its conductivity, and $I_w$ its winding current.
The implementation uses
\[
 J_z=\frac{NI_w}{S},\qquad
 R_{\rm leg}=\frac{N^2\ell}{\sigma fS}.
\]
Every active piece of the same leg shares this $J_z$. A completely overwritten
leg is an error. Resistances add within each winding. Winding resistances add
for a series circuit; conductances add for a parallel circuit. Positive
conductivity and length are required for every parallel or voltage-driven
winding. A current-driven series circuit can be solved without them, but its
resistance, voltage and power are then unavailable.

The complete example below sweeps current in a closed winding on an anisotropic
core and plots flux linkage using the existing PGFPlots snapshot interface.
\luafemmexample{static-circuits}

\subsection{Measuring a circuit and interpreting inductance}
\begin{key}{/femm/measure/quantity=circuit}
Resolve a named circuit after solving and store an independent result snapshot.
Use |\femmmeasure[name=coil-result,quantity=circuit,circuit=excitation]|.
The ordinary |parameter| key supplies a PGFPlots sweep abscissa.
\end{key}
\begin{key}{/femm/measure/circuit=\meta{name} (initially empty)}
Circuit identifier required by |quantity=circuit|. Region selections and profiles
do not select circuit legs: their circuit/winding associations do.
\end{key}

The scalar result components are |current| (A), |voltage| (V), |resistance|
($\Omega$), |power| (W), |linkage| (Wb) and |inductance| (H), subject to the
availability conditions below. Use |\femmresult|, |\femmvalue| and
|\femmresultplot| exactly as for other measurements. For magnetic depth $d$,
\[
 \lambda_w=d\sum_{\rm legs}\frac{N}{S}\int_S A_z\,dS.
\]
The triangle integral is exact for the P1 potential. Linkage is reported only
when the signed turns sum to zero in \emph{each} winding, so adding a constant
to $A_z$ cannot change it. An isolated conductor leg has no defined return path
and is not reported as a gauge-independent coil inductance.

For series windings the linkages add. For parallel windings the terminal linkage
uses the fixed DC conductance weights $R_{\rm eq}/R_w$, consistently with the
current split and magnetic energy. Problem depth is required for these measures.
At nonzero current, |inductance| is the secant ratio $\lambda/I$, and is omitted
when the model contains a permanent magnet. With multiple excited circuits it
includes their contribution to linkage; it is not an independently identified
self-inductance. Section~\ref{sec:operating} adds differential responses about a
saved bias and explains how to obtain self and mutual inductances by exciting
one port at a time. There is no automatic matrix-valued result. The resistance model includes only the declared conductor
lengths and conductivities; temperature dependence and AC losses remain separate.

\subsection{Interchange, caching and Lua entry points}
The FEM exporter writes diagonal permeability in file axes, the resolved DC
windings, and angular magnetisation laws. A rotated tensor with a nonzero cross
term cannot be represented by FEMM's two permeability entries and is rejected.
Linear sheet mixtures export their effective tensor. Parallel/voltage circuits
export series winding records at their resolved operating currents; they do not
retain a terminal voltage controller. Fill and electrical length are represented
by an effective DC conductivity. See Section~\ref{sec:interchange} for the full
contract and the distinction between source geometry and answer snapshots.

Periodic chains are split into pairs of straight primitives, each with one FEMM
periodic property. The importer accepts such straight pairs, diagonal linear
anisotropy, linear directional laminations and the restricted passive angular
expression grammar. Periodic circular-arc import, arbitrary executable angle
expressions and FEMM parallel-circuit import remain explicitly unsupported.

For spatial magnetisation, ANS output preserves the solved quadrature-averaged
source with an elementwise block representation. It is a field snapshot: use
the separate FEM export to edit the original angle law. ANS files also include
periodic nodal records. The native cache retains the original model semantics;
its revision is advanced so an older numerical cache cannot masquerade as a
solution of the extended model.

The main Lua entry points, independent of TeX, are:
\begin{itemize}
\item |luafemm-frontiers.domain(m, contours, fill_rule)|;
\item |luafemm-frontiers.path(m, name, points, closed, options)|, with
  |type|, |value|, |dx|, |dy|, |coefficient| and open-boundary parameters;
\item |luafemm-frontiers.pair(m, first, second, options)|, with
  |angle|, |x|, |y|, |cx|, |cy|, |sign| and |tolerance|;
\item |luafemm.material(m, name, properties)| additionally accepts
  |mur_x|, |mur_y|, |material_angle|, |lamination|, |lamination_fill| and
  |conductivity|;
\item |luafemm-static.region(region, options)| configures the returned region
  with |magnetization|, |center_x|, |center_y|, |order|, |circuit|, |winding|,
  |turns|, |fill| and |length|;
\item |luafemm-static.circuit(m, name, options)| and
  |luafemm-static.result(m, name)| declare and read a circuit.
\end{itemize}
Lua numerical arguments are numbers, not PGF expressions. Constructors copy
geometry and validate physical inputs. These modules are original Lua
implementations; neither a native mesher nor a native FEMM solver is needed
when compiling a document.
